\paragraph{Definition (Sidon set).}
A finite set $A\subset \mathbb{N}$ is a \emph{Sidon set} (a $B_2$ set) if
\[
a+b=c+d \ \text{with}\ a,b,c,d\in A \ \Longrightarrow\ \{a,b\}=\{c,d\}.
\]
Equivalently, the map $\{a,b\}\mapsto a+b$ from unordered pairs in $A$ is injective.

\paragraph{Definition (sumset and isolated sums).}
Write
\[
A+A := \left\{a+b:\ a,b\in A\right\}\subset \mathbb{Z}.
\]
An element $s\in A+A$ is called \emph{$\pm 1$--isolated} if
\[
s-1\notin A+A \quad\text{and}\quad s+1\notin A+A.
\]
Let
\[
I(A):=\left\lvert \left\{s\in A+A:\ s\pm 1\notin A+A\right\}\right\rvert
\]
denote the number of isolated sums.

\paragraph{Problem statement (one clean formulation).}
Is it true that
\[
\forall M\in \mathbb{N}\ \exists n_0(M)\ \forall\text{ finite Sidon }A\subset \mathbb{N}\ (\,\left\lvert A\right\rvert\ge n_0(M)\ \Rightarrow\ I(A)\ge M\,)?
\]
Equivalently, does $I(A)\to \infty$ along \emph{every} sequence of finite Sidon sets with $\left\lvert A\right\rvert\to\infty$?

\paragraph{Edge cases.}
\begin{itemize}[leftmargin=2em]
\item If $\left\lvert A\right\rvert=1$, then $A+A$ is a singleton and that unique sum is isolated.
\item If $\gcd(A)\ge 2$ (e.g.\ $A\subset 2\mathbb{Z}$), then $A+A\subset 2\mathbb{Z}$ and \emph{every} element of $A+A$ is isolated (since neighbors have opposite parity). This makes the problem nontrivial mainly in the ``$\gcd(A)=1$'' regime.
\item Translating $A$ by an integer $t$ replaces $A+A$ by $(A+A)+2t$, which preserves adjacency; hence $I(A)$ is translation invariant.
\end{itemize}
